\chapter{Mountain sickness}
\index{Mountain sickness}

\section*{Definition and Overview}
Mountain sickness, clinically referred to as acute mountain sickness (AMS) or altitude illness, is a pathological syndrome caused by rapid ascent to high altitudes (typically above 2,500 metres) without adequate acclimatisation. The condition is triggered by the physiological effects of hypobaric hypoxia, where the decrease in atmospheric pressure at high altitudes reduces the partial pressure of oxygen in the inspired air. AMS presents as a cluster of symptoms including headache, fatigue, dizziness, and gastrointestinal distress. If left unmanaged or if further ascent is attempted, AMS can progress to life-threatening complications: High Altitude Pulmonary Oedema (HAPE) or High Altitude Cerebral Oedema (HACE).

\section*{Epidemiology}
Mountain sickness is a common condition among mountain climbers, trekkers, skiers, and tourists who travel to high-altitude destinations. The incidence is directly proportional to the speed of ascent and the absolute altitude reached. Approximately 25\% of individuals ascending rapidly to 3,000 metres will develop AMS, and this rate increases to over 50\% for ascents above 4,000 metres. The condition affects individuals regardless of physical fitness, age, or sex. While mountain sickness is not a common issue within the affected region due to its relatively low-altitude ranges, it frequently affects local travellers visiting destinations such as the Andes, the Himalayas, the Rocky Mountains, or Mount Kilimanjaro.

\section*{Aetiology and Pathophysiology}
The primary driver of mountain sickness is hypobaric hypoxia, which triggers a cascade of compensatory and pathological physiological responses:
\begin{itemize}
    \item \textbf{Hypobaric hypoxia:} A decrease in barometric pressure reduces the number of oxygen molecules per breath, causing a reduction in arterial oxygen saturation (hypoxaemia).
    \item \textbf{Impaired acclimatisation:} Rapid ascent does not allow time for the body's natural acclimatisation mechanisms (such as hyperventilation, renal bicarbonate excretion to correct respiratory alkalosis, and erythropoietin-stimulated red blood cell production).
    \item \textbf{Cerebral vasodilation:} Hypoxia-induced cerebral vasodilation, combined with altered blood-brain barrier permeability, increases intracranial pressure, causing the characteristic altitude headache.
    \item \textbf{Pulmonary vasoconstriction (HAPE):} Severe hypoxia induces uneven pulmonary vasoconstriction, leading to pulmonary hypertension, elevated capillary pressure, and leakage of fluid into the alveoli (pulmonary oedema).
    \item \textbf{Brain swelling (HACE):} Severe hypoxaemia causes progressive vasogenic cerebral oedema, resulting in elevated intracranial pressure, ataxia, and altered consciousness.
\end{itemize}

\section*{Clinical Features}
The symptoms of mountain sickness typically appear within 6 to 24 hours of ascent and are classified into acute mountain sickness (AMS), HAPE, and HACE:
\begin{itemize}
    \item \textbf{Acute mountain sickness (AMS):} A throbbing headache (the hallmark symptom), accompanied by at least one of the following: loss of appetite, nausea, vomiting, dizziness, lightheadedness, fatigue, weakness, and insomnia.
    \item \textbf{High Altitude Pulmonary Oedema (HAPE) signs:} Extreme breathlessness on exertion progressing to dyspnoea at rest, a dry cough progressing to a wet cough with pink, frothy sputum, chest tightness, marked tachycardia, tachypnoea, and central cyanosis.
    \item \textbf{High Altitude Cerebral Oedema (HACE) signs:} Severe, unremitting headache, ataxia (clumsiness, loss of coordination, or a staggering "drunken" gait), confusion, hallucinations, slurred speech, and progressive lethargy.
\end{itemize}

\section*{Differential Diagnosis}
Clinicians and travel medicine specialists must distinguish mountain sickness from other illnesses common in wilderness and travel settings:
\begin{itemize}
    \item \textbf{Dehydration:} Can present with headache, fatigue, and dizziness, but is distinguished by rapid clinical improvement following oral rehydration and the absence of altitude-specific progression.
    \item \textbf{Carbon monoxide poisoning:} Often caused by poorly ventilated stoves inside tents or mountain huts, presenting with headache, nausea, and confusion.
    \item \textbf{Hypothermia:} Presents with shivering, confusion, and slurred speech, but is differentiated by measuring a low core body temperature.
    \item \textbf{Viral gastroenteritis:} Presenting with vomiting, diarrhoea, and abdominal cramps, which is usually accompanied by a fever and lacks the temporal association with altitude changes.
    \item \textbf{Migraine or tension headache:} Primary headache disorders that may be triggered by stress, travel, or change in routine, but are not accompanied by the systemic features of hypoxaemia.
\end{itemize}

\section*{When to Seek Medical Care}
Travellers should seek medical evaluation if symptoms of mild mountain sickness do not resolve with 24 hours of rest, or if they worsen despite treatment.

\textbf{Contact your local emergency services for:}
\begin{itemize}
    \item \textbf{Signs of HACE:} Ataxia (staggering gait, inability to walk a straight line), confusion, slurred speech, or extreme drowsiness.
    \item \textbf{Signs of HAPE:} Severe shortness of breath at rest, a persistent cough yielding frothy sputum, or blue-coloured lips and fingernails.
    \item \textbf{Loss of consciousness:} Unconsciousness, seizures, or inability to wake a person at high altitude.
\end{itemize}

\section*{Diagnosis and Investigations}
Diagnosis in wilderness settings is primarily clinical, supported by standardized scoring systems:
\begin{itemize}
    \item \textbf{Lake Louise Score (LLS):} A validated clinical scoring system diagnosing AMS based on the presence of a headache plus a total score of 3 or more from symptoms of fatigue, dizziness, gastrointestinal distress, and sleep disturbance.
    \item \textbf{Wilderness neurological assessment:} Assessing coordination using the tandem gait test (walking heel-to-toe in a straight line); failure indicates ataxia, confirming HACE.
    \item \textbf{Pulse oximetry:} Measuring oxygen saturation (SpO2), which will show values significantly below normal sea-level values, helping to assess the severity of hypoxia.
    \item \textbf{Chest auscultation and imaging:} Listening for crackles in the lungs; in hospital settings, chest radiographs or lung ultrasound show bilateral patchy infiltrates characteristic of HAPE.
\end{itemize}

\section*{Management}
The primary management principles for all altitude illnesses are halting ascent, resting, and descending if symptoms do not resolve:
\begin{itemize}
    \item \textbf{Immediate descent:} The definitive and most effective treatment; immediate descent of at least 500 to 1,000 metres is mandatory for anyone showing signs of HAPE or HACE.
    \item \textbf{Oxygen therapy:} Administering high-flow supplemental oxygen (2 to 4 litres per minute via nasal cannulae or face mask) to raise oxygen saturation.
    \item \textbf{Pharmacotherapy:} Acetazolamide (a carbonic anhydrase inhibitor that induces metabolic acidosis, stimulating ventilation and accelerating acclimatisation); dexamethasone (a potent corticosteroid used to treat moderate-to-severe AMS and HACE); and nifedipine (to treat HAPE by decreasing pulmonary artery pressure).
    \item \textbf{Hyperbaric chamber therapy:} Utilising a portable hyperbaric chamber (Gamow bag) in remote locations where immediate descent is impossible; inflating the chamber simulates descent by increasing barometric pressure.
    \item \textbf{Supportive care:} Halting ascent, taking simple analgesics (such as ibuprofen or paracetamol) for headaches, and maintaining hydration.
\end{itemize}

\section*{Complications}
Failure to recognise and treat mountain sickness can lead to rapid deterioration and fatal complications:
\begin{itemize}
    \item \textbf{Brain herniation (from HACE):} Progressive cerebral oedema leading to brain compression, coma, and death.
    \item \textbf{Asphyxiation (from HAPE):} Severe pulmonary oedema causing acute respiratory failure and death.
    \item \textbf{Thromboembolic events:} Increased haematocrit (polycythaemia) combined with dehydration at high altitude increases blood viscosity, raising the risk of deep vein thrombosis and pulmonary embolism.
    \item \textbf{High-altitude retinal haemorrhages:} Small haemorrhages in the retina, which are usually benign and resolve spontaneously but indicate severe systemic hypoxia.
\end{itemize}

\section*{Prognosis and Outcomes}
The prognosis for mild-to-moderate acute mountain sickness is excellent, with symptoms typically resolving completely within 24 to 48 hours of rest at the same altitude, or immediately upon descent. For patients with HACE or HAPE, the prognosis is also excellent if descent is initiated promptly and oxygen is administered. If HACE or HAPE are unrecognized or ignored, and the individual continues to ascend, the mortality rate exceeds 50\%, and death can occur within 12 to 24 hours.

\section*{Prevention and Self-Care}
Most cases of mountain sickness can be prevented through planned, gradual ascent and appropriate self-care:
\begin{itemize}
    \item \textbf{Gradual ascent rate:} Above 3,000 metres, limit the sleeping elevation gain to no more than 300 to 500 metres per day, and take a rest day every 3 to 4 days.
    \item \textbf{Climb high, sleep low:} Ascending to a higher elevation during day hikes but returning to a lower elevation to sleep to promote acclimatisation.
    \item \textbf{Pharmacological prophylaxis:} Taking acetazolamide (125 mg twice daily) starting 24 hours before ascending into high altitudes, as prescribed by a travel health specialist.
    \item \textbf{Avoid alcohol and respiratory depressants:} Refraining from alcohol, sleeping pills, and sedatives, which depress ventilation and worsen hypoxia at night.
    \item \textbf{Dietary adjustments:} Consuming a high-carbohydrate diet (which requires less oxygen to metabolise) and drinking 3 to 4 litres of water daily to prevent dehydration.
\end{itemize}

\section*{Sources and Further Reading}
The following organisations provide reliable, evidence-based information and guidelines on mountain sickness:
\begin{itemize}
    \item Travel Medicine Alliance Australia. \textit{Altitude Sickness.} https://www.travelmedicine.com.au/health-information/altitude-sickness/
    \item Healthdirect Australia. \textit{Altitude sickness.} https://www.healthdirect.gov.au/altitude-sickness
    \item NHS. \textit{Altitude sickness.} https://www.nhs.uk/conditions/altitude-sickness/
    \item Mayo Clinic. \textit{Altitude sickness: Prevention and treatment.} https://www.mayoclinic.org/
    \item Wilderness Medical Society (WMS). \textit{Clinical Practice Guidelines for the Prevention and Treatment of Acute Altitude Illness.} https://wms.org/
    \item Centers for Disease Control and Prevention (CDC). \textit{Altitude Illness.} https://wwwnc.cdc.gov/travel/yellowbook/2024/environmental-hazards-risks/altitude-illness
\end{itemize}
